% !TeX program = lualatex
% ================================================================
% Polish Tier 3 Vocabulary Version of FractionsDecimalsAndPercentagesConversions
%
% Generated by the server using the following settings:
%
%   language            Polish (pol)
%   text direction      left to right
%   font                Noto Serif
%   fallback font       Noto Serif
%   built from          latex/worksheets/percentages/fractionsDecimalsAndPercentagesConversions.tex
%   vocabulary key      latex/worksheets/percentages/fractionsDecimalsAndPercentagesConversions/L2/pol/fractionsDecimalsAndPercentagesConversions_polishKey.tex
%   tier 3 vocabulary   green   RGB 0 128 0
%   English text        black   RGB 0 0 0
%   Polish text         purple  RGB 112 48 160
%   layout macros       version 3
%
% Please don't edit any information in this comment block - the server
% compares it against the current settings and rebuilds this file when
% they differ, so changes here would be overwritten. However, feel free
% to make updates to the LaTeX below if you want to edit it.
% ================================================================

\documentclass[12pt]{article}
% ================================================================
% Copied from the English sheet this was translated from, so that
% anything it sets up for itself still works here
% ================================================================
\usepackage[a4paper,margin=1in]{geometry}
\usepackage{amsmath}
\usepackage{amssymb}
\usepackage{siunitx}
\usepackage{enumitem}
\usepackage{multicol}

% Characters this language's font has not got are borrowed from another,
% rather than being silently dropped - see docs/translated-sheet-latex.md
\directlua{%
  if luaotfload and luaotfload.add_fallback then
    luaotfload.add_fallback("eallatinfallback",
      { "Noto Serif:mode=harf;" })
  else
    tex.error("this luaotfload is too old to register a fallback font")
  end
}
\usepackage[english]{babel}
\babelprovide[import]{polish}
\babelfont[polish]{rm}[Renderer=Harfbuzz, BoldFont={Noto Serif}, ItalicFont={Noto Serif}, BoldItalicFont={Noto Serif}, RawFeature={fallback=eallatinfallback}]{Noto Serif}
\babelfont[polish]{sf}[Renderer=Harfbuzz, BoldFont={Noto Serif}, ItalicFont={Noto Serif}, BoldItalicFont={Noto Serif}, RawFeature={fallback=eallatinfallback}]{Noto Serif}
\newcommand{\ealtext}[1]{\foreignlanguage{polish}{#1}}
\newcommand{\ealtextblock}[1]{{\raggedright\foreignlanguage{polish}{#1}\par}}
% ================================================================
% EAL layout helpers
% ================================================================
\usepackage{xcolor}

\definecolor{ealkeycolour}{RGB}{0,128,0}
\definecolor{ealenglish}{RGB}{0,0,0}
\definecolor{ealtwocolour}{RGB}{112,48,160}

\newsavebox{\ealboxen}
\newsavebox{\ealboxtr}
\newlength{\ealglosswd}
\newlength{\ealglossraise}
\setlength{\ealglossraise}{1.15em}

% a tier 3 word, in English and in translation alike
\newcommand{\ealkey}[1]{\textcolor{ealkeycolour}{#1}}
\newcommand{\ealkeytr}[1]{\textcolor{ealkeycolour}{\ealtext{#1}}}

% One translated word above one English word. Built from kernel
% primitives, never a tabular, which would break wherever the sheet
% loads array, booktabs or tabularx. The box is as wide as the wider of
% the two, so a long translation pushes its neighbours aside instead of
% overlapping them, and it claims its full height so a table row or a
% TikZ node makes room for it.
\newcommand{\ealgloss}[2]{%
  \sbox{\ealboxen}{\ealkey{#1}}%
  \sbox{\ealboxtr}{\small\textcolor{ealtwocolour}{\ealtext{#2}}}%
  \setlength{\ealglosswd}{\wd\ealboxen}%
  \ifdim\wd\ealboxtr>\ealglosswd\setlength{\ealglosswd}{\wd\ealboxtr}\fi
  \leavevmode
  \makebox[\ealglosswd][c]{%
    \makebox[0pt][c]{\usebox{\ealboxen}}%
    \makebox[0pt][c]{\raisebox{\ealglossraise}{\usebox{\ealboxtr}}}%
  }%
}

% opens up line spacing around a block of glosses, so the raised
% translations sit clear of the line above
\newenvironment{ealglossed}{\par\linespread{2.1}\selectfont\ignorespaces}{\par}

% a whole translated sentence above its English counterpart. block level,
% so it does not belong inside a tabular cell
\newcommand{\ealpara}[2]{%
  \par\addvspace{0.5\baselineskip}%
  {\color{ealtwocolour}\ealtextblock{#1}}%
  \penalty200
  {\color{ealenglish}#2\par}%
  \addvspace{0.5\baselineskip}%
}

\begin{document}
\pagenumbering{gobble}

\begin{ealglossed}
\title{Conversions: \ealgloss{Decimals}{Ułamki dziesiętne}, \ealgloss{Percentages}{Procenty}, and \ealgloss{Fractions}{Ułamki}}
\date{}
\maketitle
\end{ealglossed}

\begin{enumerate}[label=\textbf{Q\arabic*.}]


\item
\begin{ealglossed}
Write the following \ealgloss{decimals}{ułamki dziesiętne} as a \ealgloss{fraction}{ułamek} in \ealgloss{simplest form}{najprostsza postać} and as a \ealgloss{percentage}{procent}:
\end{ealglossed}
\begin{multicols}{2}
\begin{enumerate}[label=(\alph*)]
    \item \( 0.125 \)
    \item \( 0.203125 \)
\end{enumerate}
\end{multicols}


\item
\begin{ealglossed}
Convert the following \ealgloss{fractions}{ułamki} into \ealgloss{decimals}{ułamki dziesiętne} and \ealgloss{percentages}{procenty}:
\end{ealglossed}
\begin{multicols}{2}
\begin{enumerate}[label=(\alph*)]
    \item \( \frac{7}{20} \)
    \item \( \frac{33}{320} \)
\end{enumerate}
\end{multicols}

\item
\begin{ealglossed}
Express the following \ealgloss{percentages}{procenty} as \ealgloss{decimals}{ułamki dziesiętne} and \ealgloss{fractions}{ułamki} in \ealgloss{simplest form}{najprostsza postać}:
\end{ealglossed}
\begin{multicols}{2}
\begin{enumerate}[label=(\alph*)]
    \item \( 17.5\% \)
    \item \( 57.8125\% \)
\end{enumerate}
\end{multicols}

\end{enumerate}


\bigskip


\begin{enumerate}[label=\textbf{Q\arabic*.}, start=4]

\item

\begin{ealglossed}
\textbf{\ealgloss{Basis points}{Punkty bazowe}} are units used in finance to describe small changes of \ealgloss{percentages}{procenty}. A \ealgloss{basis point}{punkt bazowy} is a \ealgloss{percent}{procent} of a \ealgloss{percent}{procent}:
\end{ealglossed}
\[
1 \text{ basis point} = 0.01\%
\]

\begin{ealglossed}
They are often used to describe changes in \ealgloss{interest}{odsetki} rates or \ealgloss{yields}{dochody}.
\end{ealglossed}

\begin{ealglossed}
For example, an increase of \( 25 \) \ealgloss{basis points}{punkty bazowe} means the \ealgloss{interest}{odsetki} rate goes up by \( 0.25\% \).
\end{ealglossed}

\medskip





\begin{enumerate}[label=(\alph*)]
    \item
    \begin{ealglossed}
    Write one \ealgloss{basis point}{punkt bazowy} as a \ealgloss{decimal}{ułamek dziesiętny} and a \ealgloss{fraction}{ułamek} in \ealgloss{simplest form}{najprostsza postać}.
    \end{ealglossed}
    \item
    \begin{ealglossed}
    Write 175 \ealgloss{basis points}{punkty bazowe} as a \ealgloss{decimal}{ułamek dziesiętny} and a \ealgloss{fraction}{ułamek} in \ealgloss{simplest form}{najprostsza postać}.
    \end{ealglossed}
\end{enumerate}

\bigskip

\begin{ealglossed}
    \textbf{\ealgloss{Gilts}{Brytyjskie obligacje skarbowe}} are UK government \ealgloss{bonds}{obligacje}. When you buy a \ealgloss{gilt}{brytyjska obligacja skarbowa}, you are lending money to the government.
\end{ealglossed}

\begin{ealglossed}
    The UK government currently issues 30-year \ealgloss{gilts}{brytyjskie obligacje skarbowe} offering a simple \ealgloss{return}{zwrot} (not \ealgloss{compounding}{kapitalizacja}) of $5.355\%$ per year.
\end{ealglossed}

\begin{ealglossed}
    They use these to borrow money to fund investment. Each year they pay a \textbf{\ealgloss{coupon}{kupon}} equal to $5.355\%$ of the amount borrowed. At the end of the 30 years they pay back the original amount borrowed.
\end{ealglossed}

\begin{ealglossed}
    Other durations (1-year, 5-year, 10-year) of \ealgloss{Gilts}{brytyjskie obligacje skarbowe} are also sold.
\end{ealglossed}

\medskip

\begin{enumerate}[label=(\alph*), resume]
    \item
    \begin{ealglossed}
    Find $5.355\%$ as a \ealgloss{decimal}{ułamek dziesiętny} and a \ealgloss{fraction}{ułamek} in its \ealgloss{simplest form}{najprostsza postać}.
    \end{ealglossed}
    \item Calculate the total cost for the government to borrow £1000 over 30 years, including payback of the original amount.
    \item
    \begin{ealglossed}
    Suppose the \ealgloss{coupon}{kupon} payment of a 30 year \ealgloss{gilt}{brytyjska obligacja skarbowa} rises \( 43.5 \) \ealgloss{basis points}{punkty bazowe}. How much more does it cost the government to borrow £1000 over 30 years?
    \end{ealglossed}
    % \item Gilts offer only simple interest, how might an investor achieve compounding returns, while still only investing in Gilts?
\end{enumerate}

\end{enumerate}

\end{document}